% Options for packages loaded elsewhere
\PassOptionsToPackage{unicode}{hyperref}
\PassOptionsToPackage{hyphens}{url}
\PassOptionsToPackage{dvipsnames,svgnames,x11names}{xcolor}
%
\documentclass[
  10pt,
]{article}
\usepackage{amsmath,amssymb}
\usepackage{iftex}
\ifPDFTeX
  \usepackage[T1]{fontenc}
  \usepackage[utf8]{inputenc}
  \usepackage{textcomp} % provide euro and other symbols
\else % if luatex or xetex
  \usepackage{unicode-math} % this also loads fontspec
  \defaultfontfeatures{Scale=MatchLowercase}
  \defaultfontfeatures[\rmfamily]{Ligatures=TeX,Scale=1}
\fi
\usepackage{lmodern}
\ifPDFTeX\else
  % xetex/luatex font selection
\fi
% Use upquote if available, for straight quotes in verbatim environments
\IfFileExists{upquote.sty}{\usepackage{upquote}}{}
\IfFileExists{microtype.sty}{% use microtype if available
  \usepackage[]{microtype}
  \UseMicrotypeSet[protrusion]{basicmath} % disable protrusion for tt fonts
}{}
\makeatletter
\@ifundefined{KOMAClassName}{% if non-KOMA class
  \IfFileExists{parskip.sty}{%
    \usepackage{parskip}
  }{% else
    \setlength{\parindent}{0pt}
    \setlength{\parskip}{6pt plus 2pt minus 1pt}}
}{% if KOMA class
  \KOMAoptions{parskip=half}}
\makeatother
\usepackage{xcolor}
\usepackage[a4paper,margin=2cm]{geometry}
\usepackage{color}
\usepackage{fancyvrb}
\newcommand{\VerbBar}{|}
\newcommand{\VERB}{\Verb[commandchars=\\\{\}]}
\DefineVerbatimEnvironment{Highlighting}{Verbatim}{commandchars=\\\{\}}
% Add ',fontsize=\small' for more characters per line
\usepackage{framed}
\definecolor{shadecolor}{RGB}{248,248,248}
\newenvironment{Shaded}{\begin{snugshade}}{\end{snugshade}}
\newcommand{\AlertTok}[1]{\textcolor[rgb]{0.94,0.16,0.16}{#1}}
\newcommand{\AnnotationTok}[1]{\textcolor[rgb]{0.56,0.35,0.01}{\textbf{\textit{#1}}}}
\newcommand{\AttributeTok}[1]{\textcolor[rgb]{0.13,0.29,0.53}{#1}}
\newcommand{\BaseNTok}[1]{\textcolor[rgb]{0.00,0.00,0.81}{#1}}
\newcommand{\BuiltInTok}[1]{#1}
\newcommand{\CharTok}[1]{\textcolor[rgb]{0.31,0.60,0.02}{#1}}
\newcommand{\CommentTok}[1]{\textcolor[rgb]{0.56,0.35,0.01}{\textit{#1}}}
\newcommand{\CommentVarTok}[1]{\textcolor[rgb]{0.56,0.35,0.01}{\textbf{\textit{#1}}}}
\newcommand{\ConstantTok}[1]{\textcolor[rgb]{0.56,0.35,0.01}{#1}}
\newcommand{\ControlFlowTok}[1]{\textcolor[rgb]{0.13,0.29,0.53}{\textbf{#1}}}
\newcommand{\DataTypeTok}[1]{\textcolor[rgb]{0.13,0.29,0.53}{#1}}
\newcommand{\DecValTok}[1]{\textcolor[rgb]{0.00,0.00,0.81}{#1}}
\newcommand{\DocumentationTok}[1]{\textcolor[rgb]{0.56,0.35,0.01}{\textbf{\textit{#1}}}}
\newcommand{\ErrorTok}[1]{\textcolor[rgb]{0.64,0.00,0.00}{\textbf{#1}}}
\newcommand{\ExtensionTok}[1]{#1}
\newcommand{\FloatTok}[1]{\textcolor[rgb]{0.00,0.00,0.81}{#1}}
\newcommand{\FunctionTok}[1]{\textcolor[rgb]{0.13,0.29,0.53}{\textbf{#1}}}
\newcommand{\ImportTok}[1]{#1}
\newcommand{\InformationTok}[1]{\textcolor[rgb]{0.56,0.35,0.01}{\textbf{\textit{#1}}}}
\newcommand{\KeywordTok}[1]{\textcolor[rgb]{0.13,0.29,0.53}{\textbf{#1}}}
\newcommand{\NormalTok}[1]{#1}
\newcommand{\OperatorTok}[1]{\textcolor[rgb]{0.81,0.36,0.00}{\textbf{#1}}}
\newcommand{\OtherTok}[1]{\textcolor[rgb]{0.56,0.35,0.01}{#1}}
\newcommand{\PreprocessorTok}[1]{\textcolor[rgb]{0.56,0.35,0.01}{\textit{#1}}}
\newcommand{\RegionMarkerTok}[1]{#1}
\newcommand{\SpecialCharTok}[1]{\textcolor[rgb]{0.81,0.36,0.00}{\textbf{#1}}}
\newcommand{\SpecialStringTok}[1]{\textcolor[rgb]{0.31,0.60,0.02}{#1}}
\newcommand{\StringTok}[1]{\textcolor[rgb]{0.31,0.60,0.02}{#1}}
\newcommand{\VariableTok}[1]{\textcolor[rgb]{0.00,0.00,0.00}{#1}}
\newcommand{\VerbatimStringTok}[1]{\textcolor[rgb]{0.31,0.60,0.02}{#1}}
\newcommand{\WarningTok}[1]{\textcolor[rgb]{0.56,0.35,0.01}{\textbf{\textit{#1}}}}
\usepackage{longtable,booktabs,array}
\usepackage{calc} % for calculating minipage widths
% Correct order of tables after \paragraph or \subparagraph
\usepackage{etoolbox}
\makeatletter
\patchcmd\longtable{\par}{\if@noskipsec\mbox{}\fi\par}{}{}
\makeatother
% Allow footnotes in longtable head/foot
\IfFileExists{footnotehyper.sty}{\usepackage{footnotehyper}}{\usepackage{footnote}}
\makesavenoteenv{longtable}
\setlength{\emergencystretch}{3em} % prevent overfull lines
\providecommand{\tightlist}{%
  \setlength{\itemsep}{0pt}\setlength{\parskip}{0pt}}
\setcounter{secnumdepth}{-\maxdimen} % remove section numbering
\ifLuaTeX
\usepackage[bidi=basic]{babel}
\else
\usepackage[bidi=default]{babel}
\fi
\babelprovide[main,import]{brazilian}
% get rid of language-specific shorthands (see #6817):
\let\LanguageShortHands\languageshorthands
\def\languageshorthands#1{}
\usepackage{fvextra}
\DefineVerbatimEnvironment{Highlighting}{Verbatim}{breaklines,breakanywhere,commandchars=\\\{\}}
\usepackage{newunicodechar}
\usepackage{amssymb}
\newunicodechar{✓}{\ensuremath{\checkmark}}
\newunicodechar{⋮}{\ensuremath{\vdots}}
\newunicodechar{→}{\ensuremath{\rightarrow}}
\usepackage{xcolor}
\definecolor{velmont}{HTML}{1D4ED8}
\usepackage{sectsty}
\allsectionsfont{\color{velmont}\sffamily}
\usepackage{fancyhdr}
\pagestyle{fancy}\fancyhf{}
\lhead{\small\sffamily GestorEng · ZimaOS}\rhead{\small\sffamily Velmont}
\cfoot{\small\thepage}
\renewcommand{\headrulewidth}{0.4pt}
\let\oldquote\quote\let\endoldquote\endquote
\renewenvironment{quote}{\begin{oldquote}\small\itshape}{\end{oldquote}}
\setlength{\parskip}{0.5em}
\newunicodechar{─}{-}
\ifLuaTeX
  \usepackage{selnolig}  % disable illegal ligatures
\fi
\IfFileExists{bookmark.sty}{\usepackage{bookmark}}{\usepackage{hyperref}}
\IfFileExists{xurl.sty}{\usepackage{xurl}}{} % add URL line breaks if available
\urlstyle{same}
\hypersetup{
  pdftitle={GestorEng no ZimaOS --- Guia de Instalação},
  pdflang={pt-BR},
  colorlinks=true,
  linkcolor={blue},
  filecolor={Maroon},
  citecolor={Blue},
  urlcolor={blue},
  pdfcreator={LaTeX via pandoc}}

\title{GestorEng no ZimaOS --- Guia de Instalação}
\usepackage{etoolbox}
\makeatletter
\providecommand{\subtitle}[1]{% add subtitle to \maketitle
  \apptocmd{\@title}{\par {\large #1 \par}}{}{}
}
\makeatother
\subtitle{Velmont · banco interno ou externo · passo a passo}
\author{}
\date{16/11/2026}

\begin{document}
\maketitle

{
\hypersetup{linkcolor=}
\setcounter{tocdepth}{1}
\tableofcontents
}
\hypertarget{antes-de-comeuxe7ar}{%
\section{Antes de começar}\label{antes-de-comeuxe7ar}}

\hypertarget{qual-arquivo-usar}{%
\subsection{Qual arquivo usar}\label{qual-arquivo-usar}}

\begin{longtable}[]{@{}
  >{\raggedright\arraybackslash}p{(\columnwidth - 4\tabcolsep) * \real{0.3333}}
  >{\raggedright\arraybackslash}p{(\columnwidth - 4\tabcolsep) * \real{0.3333}}
  >{\raggedright\arraybackslash}p{(\columnwidth - 4\tabcolsep) * \real{0.3333}}@{}}
\toprule\noalign{}
\begin{minipage}[b]{\linewidth}\raggedright
Arquivo
\end{minipage} & \begin{minipage}[b]{\linewidth}\raggedright
Quando usar
\end{minipage} & \begin{minipage}[b]{\linewidth}\raggedright
O que sobe
\end{minipage} \\
\midrule\noalign{}
\endhead
\bottomrule\noalign{}
\endlastfoot
\texttt{banco-interno.yml} & Instalação nova e independente
(\textbf{recomendado} se você não tem um Postgres dedicado) & PostgreSQL
16, Redis, API, worker e web \\
\texttt{banco-externo.yml} & Você \textbf{já tem} PostgreSQL (13 a 18)
e/ou Redis rodando na rede & API, worker e web \\
\end{longtable}

Os dois estão em \texttt{deploy/zimaos/} no repositório. As variáveis
ficam \textbf{direto em cada serviço} (\texttt{environment:}), então
aparecem e podem ser alteradas no próprio painel do ZimaOS.

\hypertarget{o-que-vocuxea-vai-precisar}{%
\subsection{O que você vai precisar}\label{o-que-vocuxea-vai-precisar}}

\begin{itemize}
\tightlist
\item
  ZimaOS com acesso ao \textbf{terminal}: pelo painel (Configurações
  \textgreater{} Terminal) ou por SSH
  (\texttt{ssh\ SEU\_USUARIO@IP-DO-ZIMAOS}).
\item
  Uma conta do GitHub com acesso ao repositório
  \texttt{incorpore/GestorEng} (o repositório é \textbf{privado}).
\item
  Cerca de 2 GB livres de disco.
\end{itemize}

\begin{quote}
\textbf{Convenção deste guia:} \texttt{IP-DO-ZIMAOS} é o endereço do seu
servidor (ex.: \texttt{192.168.3.100}). Comandos que começam com
\texttt{sudo} são executados \textbf{no terminal do ZimaOS}. Comandos
com \texttt{docker\ exec} executam algo \textbf{dentro de um container}.
\end{quote}

\hypertarget{passo-1-liberar-o-download-das-imagens-reposituxf3rio-privado}{%
\section{Passo 1 --- Liberar o download das imagens (repositório
privado)}\label{passo-1-liberar-o-download-das-imagens-reposituxf3rio-privado}}

Como o repositório é privado, as imagens
\texttt{ghcr.io/incorpore/gestoreng-api} e \texttt{gestoreng-web} também
são privadas. O ZimaOS precisa de um \textbf{token do GitHub} para
baixá-las.

\textbf{1.1 Criar o token (no navegador):}

\begin{enumerate}
\def\labelenumi{\arabic{enumi}.}
\tightlist
\item
  GitHub \textgreater{} clique na sua foto \textgreater{}
  \textbf{Settings} \textgreater{} \textbf{Developer settings}
  \textgreater{} \textbf{Personal access tokens} \textgreater{}
  \textbf{Tokens (classic)}.
\item
  \textbf{Generate new token (classic)}.
\item
  \emph{Note}: \texttt{ZimaOS\ GestorEng} · \emph{Expiration}: o prazo
  que preferir (anote para renovar).
\item
  Marque \textbf{somente} \texttt{read:packages}.
\item
  \textbf{Generate token} e copie o valor (começa com \texttt{ghp\_}).
  Ele aparece uma única vez.
\end{enumerate}

\textbf{1.2 Fazer login no ZimaOS (terminal):}

\begin{Shaded}
\begin{Highlighting}[]
\BuiltInTok{echo} \StringTok{"ghp\_SEU\_TOKEN"} \KeywordTok{|} \FunctionTok{sudo}\NormalTok{ docker login ghcr.io }\AttributeTok{{-}u}\NormalTok{ SEU\_USUARIO\_GITHUB }\AttributeTok{{-}{-}password{-}stdin}
\end{Highlighting}
\end{Shaded}

Deve aparecer \texttt{Login\ Succeeded}. O login fica salvo; só repita
se o token expirar.

\textbf{1.3 Baixar as imagens:}

\begin{Shaded}
\begin{Highlighting}[]
\FunctionTok{sudo}\NormalTok{ docker pull ghcr.io/incorpore/gestoreng{-}api:latest}
\FunctionTok{sudo}\NormalTok{ docker pull ghcr.io/incorpore/gestoreng{-}web:latest}
\end{Highlighting}
\end{Shaded}

\textbf{Alternativa --- tornar somente as imagens públicas} (o
repositório e o código no GitHub continuam privados; o ZimaOS baixa sem
login, e os passos 1.1 e 1.2 deixam de ser necessários):

\begin{enumerate}
\def\labelenumi{\arabic{enumi}.}
\tightlist
\item
  GitHub \textgreater{} organização/perfil \textbf{incorpore}
  \textgreater{} aba \textbf{Packages} \textgreater{}
  \textbf{gestoreng-api}.
\item
  \textbf{Package settings} (lateral direita) \textgreater{}
  \emph{Danger Zone} \textgreater{} \textbf{Change visibility}
  \textgreater{} \textbf{Public} \textgreater{} digite o nome para
  confirmar.
\item
  Repita para \textbf{gestoreng-web}.
\end{enumerate}

Se a opção \emph{Public} estiver bloqueada (conta de organização):
\emph{Settings} da organização \textgreater{} \textbf{Packages}
\textgreater{} em \emph{Package creation}, permita pacotes
\textbf{Public}.

\begin{quote}
Com a imagem pública, qualquer pessoa pode baixá-la e ler o código do
sistema que vai dentro dela. Nenhuma senha ou chave fica na imagem (elas
ficam só no seu arquivo \texttt{.yml}), mas, se o código for sigiloso,
mantenha as imagens privadas e use o login dos passos 1.1 e 1.2.
\end{quote}

\begin{quote}
\textbf{Erro \texttt{lookup\ ghcr.1panel.live:\ no\ such\ host}?} O
Docker do servidor está usando um espelho que saiu do ar. Veja a seção
\emph{Problemas comuns} no final e repita o 1.3.
\end{quote}

\hypertarget{passo-2-criar-as-pastas-de-dados}{%
\section{Passo 2 --- Criar as pastas de
dados}\label{passo-2-criar-as-pastas-de-dados}}

\begin{Shaded}
\begin{Highlighting}[]
\FunctionTok{sudo}\NormalTok{ mkdir }\AttributeTok{{-}p}\NormalTok{ /media/URSUS/AppData/GestorEng/storage}
\FunctionTok{sudo}\NormalTok{ chown }\AttributeTok{{-}R}\NormalTok{ 10001:10001 /media/URSUS/AppData/GestorEng/storage}
\end{Highlighting}
\end{Shaded}

Somente para \textbf{banco interno}, crie também a pasta do PostgreSQL
(o próprio Postgres ajusta as permissões):

\begin{Shaded}
\begin{Highlighting}[]
\FunctionTok{sudo}\NormalTok{ mkdir }\AttributeTok{{-}p}\NormalTok{ /media/URSUS/AppData/GestorEng/postgres}
\end{Highlighting}
\end{Shaded}

\begin{quote}
Se quiser outro local (outro disco, NAS montado), troque o caminho
\texttt{/media/URSUS/AppData/GestorEng} em \textbf{todas} as linhas
\texttt{source:} do arquivo \texttt{.yml}.
\end{quote}

\hypertarget{passo-3-gerar-as-chaves-secretas}{%
\section{Passo 3 --- Gerar as chaves
secretas}\label{passo-3-gerar-as-chaves-secretas}}

Execute \textbf{duas vezes} e guarde os dois resultados em local seguro
(gerenciador de senhas):

\begin{Shaded}
\begin{Highlighting}[]
\ExtensionTok{openssl}\NormalTok{ rand }\AttributeTok{{-}hex}\NormalTok{ 32}
\end{Highlighting}
\end{Shaded}

\begin{itemize}
\tightlist
\item
  O primeiro valor será a \texttt{SECRET\_KEY}.
\item
  O segundo será a \texttt{ENCRYPTION\_KEY}.
\end{itemize}

\begin{quote}
\textbf{Atenção:} a \texttt{ENCRYPTION\_KEY} protege as chaves do
OpenRouter, Documenso e do backup S3 gravadas no banco. Se ela for
perdida ou trocada, será preciso cadastrar essas integrações novamente.
\textbf{Não troque depois.}
\end{quote}

\hypertarget{passo-4-preparar-o-banco-externo-somente-banco-externo.yml}{%
\section{\texorpdfstring{Passo 4 --- Preparar o banco externo (somente
\texttt{banco-externo.yml})}{Passo 4 --- Preparar o banco externo (somente banco-externo.yml)}}\label{passo-4-preparar-o-banco-externo-somente-banco-externo.yml}}

Pule este passo se for usar o banco interno.

\textbf{4.1 Criar usuário e banco no PostgreSQL.} Descubra o nome do
container do seu Postgres e entre nele:

\begin{Shaded}
\begin{Highlighting}[]
\FunctionTok{sudo}\NormalTok{ docker ps }\AttributeTok{{-}{-}format} \StringTok{\textquotesingle{}\{\{.Names\}\}  \{\{.Image\}\}\textquotesingle{}} \KeywordTok{|} \FunctionTok{grep} \AttributeTok{{-}i}\NormalTok{ postgres}
\FunctionTok{sudo}\NormalTok{ docker exec }\AttributeTok{{-}it}\NormalTok{ NOME\_DO\_CONTAINER\_POSTGRES psql }\AttributeTok{{-}U}\NormalTok{ postgres}
\end{Highlighting}
\end{Shaded}

Dentro do \texttt{psql}, execute (troque a senha):

\begin{Shaded}
\begin{Highlighting}[]
\KeywordTok{CREATE} \FunctionTok{USER}\NormalTok{ gestoreng }\KeywordTok{WITH} \KeywordTok{PASSWORD} \StringTok{\textquotesingle{}SUA\_SENHA\_DO\_BANCO\textquotesingle{}}\NormalTok{;}
\KeywordTok{CREATE} \KeywordTok{DATABASE}\NormalTok{ gestoreng OWNER gestoreng;}
\NormalTok{\textbackslash{}q}
\end{Highlighting}
\end{Shaded}

\begin{quote}
Use nomes \textbf{em minúsculas} (\texttt{gestoreng}). Se você escrever
\texttt{GestorEng} sem aspas, o PostgreSQL grava como \texttt{gestoreng}
e a conexão com o nome em maiúsculas falha.
\end{quote}

\textbf{4.2 Testar se o servidor responde} (a partir do ZimaOS):

\begin{Shaded}
\begin{Highlighting}[]
\FunctionTok{sudo}\NormalTok{ docker run }\AttributeTok{{-}{-}rm}\NormalTok{ postgres:16{-}alpine pg\_isready }\AttributeTok{{-}h}\NormalTok{ 192.168.3.100 }\AttributeTok{{-}p}\NormalTok{ 5432}
\CommentTok{\# esperado: 192.168.3.100:5432 {-} accepting connections}
\end{Highlighting}
\end{Shaded}

Se usar o \textbf{Redis externo}, teste também:

\begin{Shaded}
\begin{Highlighting}[]
\FunctionTok{sudo}\NormalTok{ docker run }\AttributeTok{{-}{-}rm}\NormalTok{ redis:7{-}alpine redis{-}cli }\AttributeTok{{-}h}\NormalTok{ 192.168.3.100 }\AttributeTok{{-}p}\NormalTok{ 6379 }\AttributeTok{{-}a} \StringTok{\textquotesingle{}SENHA\_DO\_REDIS\textquotesingle{}} \AttributeTok{{-}{-}no{-}auth{-}warning}\NormalTok{ ping}
\CommentTok{\# esperado: PONG}
\end{Highlighting}
\end{Shaded}

Se algum teste falhar, verifique se a porta está publicada no container
do banco (\texttt{ports:\ 5432:5432} / \texttt{6379:6379}) e se o
firewall permite a conexão.

\hypertarget{passo-5-preencher-o-arquivo-de-instalauxe7uxe3o}{%
\section{Passo 5 --- Preencher o arquivo de
instalação}\label{passo-5-preencher-o-arquivo-de-instalauxe7uxe3o}}

Abra o arquivo escolhido (\texttt{banco-interno.yml} ou
\texttt{banco-externo.yml}) e troque todos os valores
\texttt{TROQUE\_*}.

\begin{quote}
\textbf{Regra de ouro:} os serviços \textbf{\texttt{api} e
\texttt{worker} têm as mesmas variáveis e elas devem ficar iguais}
(principalmente senhas, \texttt{SECRET\_KEY} e
\texttt{ENCRYPTION\_KEY}). No banco interno, a senha do banco também
aparece no serviço \texttt{db} (\texttt{POSTGRES\_PASSWORD}) e a do
Redis aparece no \texttt{command} do serviço \texttt{redis} (linha após
\texttt{-\/-requirepass}). Dica: use \emph{Localizar e substituir} do
editor de texto para trocar cada \texttt{TROQUE\_*} em todos os lugares
de uma vez.
\end{quote}

\begin{longtable}[]{@{}
  >{\raggedright\arraybackslash}p{(\columnwidth - 2\tabcolsep) * \real{0.5000}}
  >{\raggedright\arraybackslash}p{(\columnwidth - 2\tabcolsep) * \real{0.5000}}@{}}
\toprule\noalign{}
\begin{minipage}[b]{\linewidth}\raggedright
Variável
\end{minipage} & \begin{minipage}[b]{\linewidth}\raggedright
O que colocar
\end{minipage} \\
\midrule\noalign{}
\endhead
\bottomrule\noalign{}
\endlastfoot
\texttt{POSTGRES\_PASSWORD} & Senha do banco (interno: será criada
agora; externo: a do passo 4.1) \\
\texttt{POSTGRES\_HOST}, \texttt{POSTGRES\_DB}, \texttt{POSTGRES\_USER}
& Só no externo: endereço, banco e usuário do passo 4 \\
\texttt{REDIS\_PASSWORD} & Senha do Redis (interno: escolha uma e repita
no \texttt{command} do serviço \texttt{redis}; externo: a existente) \\
\texttt{REDIS\_PREFIX} & Prefixo das chaves no Redis (ex.:
\texttt{GestorEng\_}) \\
\texttt{SECRET\_KEY} / \texttt{ENCRYPTION\_KEY} & Os dois valores do
passo 3 \\
\texttt{ADMIN\_EMAIL} / \texttt{ADMIN\_PASSWORD} & Seu acesso inicial
(10+ caracteres, misturando maiúsculas, minúsculas, números/símbolos) \\
\texttt{COOKIE\_SECURE} & \texttt{"false"} acessando por
\texttt{http://IP:8917}; \texttt{"true"} se usar HTTPS \\
\texttt{CORS\_ORIGINS} & Troque \texttt{192.168.3.100} pelo IP do ZimaOS
(e acrescente seu domínio, se tiver) \\
\texttt{MAX\_UPLOAD\_MB} & Tamanho máximo de upload; mantenha
\texttt{MAX\_BODY} do serviço \texttt{web} igual (ex.: \texttt{50} e
\texttt{50m}) \\
Porta \texttt{published:\ "8917"} (serviço \texttt{web}) & Porta de
acesso; mantenha igual a \texttt{port\_map} no final do arquivo \\
\end{longtable}

As senhas podem conter \textbf{qualquer caractere} --- o sistema monta
as URLs de conexão e codifica tudo sozinho.

\hypertarget{passo-6-instalar-no-zimaos}{%
\section{Passo 6 --- Instalar no
ZimaOS}\label{passo-6-instalar-no-zimaos}}

\textbf{Opção A --- pelo painel (recomendado):}

\begin{enumerate}
\def\labelenumi{\arabic{enumi}.}
\tightlist
\item
  Abra o painel do ZimaOS \textgreater{} \textbf{App Store}.
\item
  Clique em \textbf{+} (canto superior) \textgreater{} \textbf{Install a
  customized app} (\emph{Instalar app personalizado}).
\item
  Clique em \textbf{Import} (ícone no canto superior direito da janela)
  \textgreater{} aba \textbf{Docker Compose}.
\item
  Cole o conteúdo \textbf{inteiro} do arquivo \texttt{.yml} já
  preenchido \textgreater{} \textbf{Submit}.
\item
  Confira os dados exibidos e clique em \textbf{Install}.
\end{enumerate}

\textbf{Opção B --- pelo terminal:}

\begin{Shaded}
\begin{Highlighting}[]
\FunctionTok{sudo}\NormalTok{ mkdir }\AttributeTok{{-}p}\NormalTok{ /media/URSUS/AppData/GestorEng/compose}
\FunctionTok{sudo}\NormalTok{ nano /media/URSUS/AppData/GestorEng/compose/docker{-}compose.yml   }\CommentTok{\# cole o .yml preenchido, Ctrl+O, Enter, Ctrl+X}
\BuiltInTok{cd}\NormalTok{ /media/URSUS/AppData/GestorEng/compose}
\FunctionTok{sudo}\NormalTok{ docker compose up }\AttributeTok{{-}d}
\end{Highlighting}
\end{Shaded}

\begin{quote}
Na opção B o app não aparece com ícone no painel da mesma forma; prefira
a opção A.
\end{quote}

\hypertarget{passo-7-conferir-se-estuxe1-tudo-funcionando}{%
\section{Passo 7 --- Conferir se está tudo
funcionando}\label{passo-7-conferir-se-estuxe1-tudo-funcionando}}

\textbf{7.1 Acompanhar a inicialização da API} (a primeira vez leva até
1 minuto, pois cria as tabelas):

\begin{Shaded}
\begin{Highlighting}[]
\FunctionTok{sudo}\NormalTok{ docker logs }\AttributeTok{{-}f}\NormalTok{ gestoreng{-}api}
\CommentTok{\# aguarde: "Application startup complete" (Ctrl+C para sair)}
\end{Highlighting}
\end{Shaded}

\textbf{7.2 Diagnóstico completo} (banco, Redis, pasta de arquivos e
backup):

\begin{Shaded}
\begin{Highlighting}[]
\FunctionTok{sudo}\NormalTok{ docker exec gestoreng{-}api python }\AttributeTok{{-}m}\NormalTok{ app.manage verificar}
\end{Highlighting}
\end{Shaded}

Saída esperada: linhas com \texttt{✓} para PostgreSQL, Redis, pasta de
arquivos e \texttt{pg\_dump}.

\textbf{7.3 Acessar:} \texttt{http://IP-DO-ZIMAOS:8917} e entrar com
\texttt{ADMIN\_EMAIL} / \texttt{ADMIN\_PASSWORD}.

\textbf{7.4 Primeiras configurações no sistema:}

\begin{enumerate}
\def\labelenumi{\arabic{enumi}.}
\tightlist
\item
  \textbf{Minha conta} (canto superior direito) \textgreater{} ative a
  \textbf{verificação em duas etapas}.
\item
  \textbf{Configurações \textgreater{} Empresa}: dados, logo e cor da
  marca.
\item
  \textbf{Configurações \textgreater{} Papel timbrado}: envie o
  PDF/imagem A4 e ajuste as margens.
\item
  \textbf{Configurações \textgreater{} Integrações}: OpenRouter e
  Documenso.
\item
  \textbf{Configurações \textgreater{} Backup}: bucket S3 e horário
  (fortemente recomendado).
\item
  \textbf{Usuários e acesso}: crie os usuários e atribua perfis.
\end{enumerate}

\hypertarget{comandos-de-manutenuxe7uxe3o-terminal-do-container}{%
\section{Comandos de manutenção (terminal do
container)}\label{comandos-de-manutenuxe7uxe3o-terminal-do-container}}

Todos são executados a partir do terminal do ZimaOS com
\texttt{sudo\ docker\ exec\ -it\ gestoreng-api\ ...}:

\begin{longtable}[]{@{}
  >{\raggedright\arraybackslash}p{(\columnwidth - 2\tabcolsep) * \real{0.5000}}
  >{\raggedright\arraybackslash}p{(\columnwidth - 2\tabcolsep) * \real{0.5000}}@{}}
\toprule\noalign{}
\begin{minipage}[b]{\linewidth}\raggedright
Para\ldots{}
\end{minipage} & \begin{minipage}[b]{\linewidth}\raggedright
Comando
\end{minipage} \\
\midrule\noalign{}
\endhead
\bottomrule\noalign{}
\endlastfoot
Ver a ajuda &
\texttt{sudo\ docker\ exec\ -it\ gestoreng-api\ python\ -m\ app.manage} \\
Listar usuários &
\texttt{sudo\ docker\ exec\ -it\ gestoreng-api\ python\ -m\ app.manage\ usuarios} \\
Criar outro administrador &
\texttt{sudo\ docker\ exec\ -it\ gestoreng-api\ python\ -m\ app.manage\ criar-admin\ email@velmont.com.br\ Nome\ Completo} \\
Redefinir senha (esqueci a senha) &
\texttt{sudo\ docker\ exec\ -it\ gestoreng-api\ python\ -m\ app.manage\ senha\ email@velmont.com.br} \\
Desbloquear usuário &
\texttt{sudo\ docker\ exec\ -it\ gestoreng-api\ python\ -m\ app.manage\ desbloquear\ email@velmont.com.br} \\
Remover 2FA (perdeu o celular) &
\texttt{sudo\ docker\ exec\ -it\ gestoreng-api\ python\ -m\ app.manage\ remover-2fa\ email@velmont.com.br} \\
Encerrar sessões de um usuário &
\texttt{sudo\ docker\ exec\ -it\ gestoreng-api\ python\ -m\ app.manage\ sair-todos\ email@velmont.com.br} \\
Diagnóstico &
\texttt{sudo\ docker\ exec\ -it\ gestoreng-api\ python\ -m\ app.manage\ verificar} \\
Gerar nova chave &
\texttt{sudo\ docker\ exec\ -it\ gestoreng-api\ python\ -m\ app.manage\ gerar-chave} \\
Executar backup agora &
\texttt{sudo\ docker\ exec\ -it\ gestoreng-api\ ./docker-entrypoint.sh\ backup\ executar} \\
Listar backups no S3 &
\texttt{sudo\ docker\ exec\ -it\ gestoreng-api\ ./docker-entrypoint.sh\ backup\ listar} \\
\end{longtable}

Os comandos que pedem senha a solicitam de forma oculta (não aparece no
histórico do terminal).

\textbf{Acesso direto ao banco interno} (consultas manuais):

\begin{Shaded}
\begin{Highlighting}[]
\FunctionTok{sudo}\NormalTok{ docker exec }\AttributeTok{{-}it}\NormalTok{ gestoreng{-}db psql }\AttributeTok{{-}U}\NormalTok{ gestoreng }\AttributeTok{{-}d}\NormalTok{ gestoreng}
\end{Highlighting}
\end{Shaded}

\hypertarget{alterar-variuxe1veis-depois-da-instalauxe7uxe3o}{%
\section{Alterar variáveis depois da
instalação}\label{alterar-variuxe1veis-depois-da-instalauxe7uxe3o}}

\textbf{Pelo painel:} ícone do GestorEng \textgreater{} menu (⋮)
\textgreater{} \textbf{Settings} (\emph{Configurações}). Cada container
(\texttt{api}, \texttt{worker}, \texttt{web}, \texttt{db},
\texttt{redis}) tem sua aba com as \textbf{variáveis de ambiente} ---
altere o valor e clique em \textbf{Save}; o ZimaOS recria o container.
\textbf{Ao mudar algo na \texttt{api}, faça a mesma alteração no
\texttt{worker}.}

\textbf{Pelo terminal (se instalou pela opção B):} edite o arquivo e
rode \texttt{sudo\ docker\ compose\ up\ -d} na pasta dele.

Efeito de cada alteração:

\begin{longtable}[]{@{}
  >{\raggedright\arraybackslash}p{(\columnwidth - 4\tabcolsep) * \real{0.3333}}
  >{\raggedright\arraybackslash}p{(\columnwidth - 4\tabcolsep) * \real{0.3333}}
  >{\raggedright\arraybackslash}p{(\columnwidth - 4\tabcolsep) * \real{0.3333}}@{}}
\toprule\noalign{}
\begin{minipage}[b]{\linewidth}\raggedright
Variável
\end{minipage} & \begin{minipage}[b]{\linewidth}\raggedright
Pode trocar?
\end{minipage} & \begin{minipage}[b]{\linewidth}\raggedright
Observação
\end{minipage} \\
\midrule\noalign{}
\endhead
\bottomrule\noalign{}
\endlastfoot
\texttt{COOKIE\_SECURE}, \texttt{CORS\_ORIGINS}, \texttt{TZ},
\texttt{WEB\_CONCURRENCY}, \texttt{REDIS\_PREFIX} & Sim & Basta
salvar. \\
\texttt{MAX\_UPLOAD\_MB} + \texttt{MAX\_BODY} & Sim & Altere \textbf{os
dois} juntos. \\
Porta (\texttt{published:\ "8917"}) & Sim & Troque também
\texttt{port\_map} no fim do arquivo e o IP:porta em
\texttt{CORS\_ORIGINS}. \\
\texttt{SECRET\_KEY} & Sim & Todos os usuários serão desconectados
(fazem login de novo). \\
\texttt{ENCRYPTION\_KEY} & \textbf{Não} & Se trocar, recadastre
OpenRouter, Documenso e o backup S3. \\
\texttt{ADMIN\_EMAIL} / \texttt{ADMIN\_PASSWORD} & Sem efeito & Só valem
na 1ª instalação. Use os comandos \texttt{criar-admin} /
\texttt{senha}. \\
Senha do Redis (interno) & Sim & Troque no \texttt{command} do
\texttt{redis} e em \texttt{REDIS\_PASSWORD} de \texttt{api} e
\texttt{worker}. \\
\texttt{POSTGRES\_PASSWORD} (interno) & Sim, em 2 etapas & Veja
abaixo. \\
\texttt{POSTGRES\_PASSWORD} (externo) & Sim, em 2 etapas & Troque
primeiro no seu PostgreSQL, depois no arquivo. \\
Pastas (\texttt{source:}) & Com cuidado & Pare o app, \textbf{copie} os
dados para a nova pasta, ajuste o \texttt{chown}, altere o arquivo. \\
\end{longtable}

\textbf{Trocar a senha do banco interno} (a variável só é lida na
criação do banco):

\begin{Shaded}
\begin{Highlighting}[]
\FunctionTok{sudo}\NormalTok{ docker exec }\AttributeTok{{-}it}\NormalTok{ gestoreng{-}db psql }\AttributeTok{{-}U}\NormalTok{ gestoreng }\AttributeTok{{-}d}\NormalTok{ gestoreng }\AttributeTok{{-}c} \StringTok{"ALTER USER gestoreng PASSWORD \textquotesingle{}NOVA\_SENHA\textquotesingle{};"}
\end{Highlighting}
\end{Shaded}

Depois altere \texttt{POSTGRES\_PASSWORD} para \texttt{NOVA\_SENHA} nos
serviços \texttt{db}, \texttt{api} e \texttt{worker} e salve.

\hypertarget{atualizar-o-gestoreng-para-uma-nova-versuxe3o}{%
\section{Atualizar o GestorEng para uma nova
versão}\label{atualizar-o-gestoreng-para-uma-nova-versuxe3o}}

\begin{Shaded}
\begin{Highlighting}[]
\FunctionTok{sudo}\NormalTok{ docker pull ghcr.io/incorpore/gestoreng{-}api:latest}
\FunctionTok{sudo}\NormalTok{ docker pull ghcr.io/incorpore/gestoreng{-}web:latest}
\end{Highlighting}
\end{Shaded}

Depois, no painel: ícone do GestorEng \textgreater{} \textbf{Restart}
não basta --- use \textbf{Settings \textgreater{} Save} (recria os
containers com as imagens novas). Pelo terminal:
\texttt{sudo\ docker\ compose\ up\ -d} na pasta do arquivo. As
atualizações do banco (migrações) são aplicadas automaticamente na
inicialização da API. Faça um backup antes de atualizar.

\hypertarget{backup-e-restaurauxe7uxe3o}{%
\section{Backup e restauração}\label{backup-e-restaurauxe7uxe3o}}

\textbf{Backup automático (recomendado):} configure em
\emph{Configurações \textgreater{} Backup} (S3, MinIO, Backblaze,
Cloudflare R2\ldots). Inclui banco e arquivos, roda todo dia no horário
escolhido e apaga os antigos conforme a retenção.

\textbf{Backup manual extra do banco interno} (arquivo na pasta do app):

\begin{Shaded}
\begin{Highlighting}[]
\FunctionTok{sudo}\NormalTok{ docker exec gestoreng{-}db pg\_dump }\AttributeTok{{-}U}\NormalTok{ gestoreng }\AttributeTok{{-}Fc}\NormalTok{ gestoreng }\KeywordTok{|} \FunctionTok{sudo}\NormalTok{ tee /media/URSUS/AppData/GestorEng/manual{-}}\VariableTok{$(}\FunctionTok{date}\NormalTok{ +\%F}\VariableTok{)}\NormalTok{.dump }\OperatorTok{\textgreater{}}\NormalTok{ /dev/null}
\end{Highlighting}
\end{Shaded}

\textbf{Restaurar um backup do S3} (sobrescreve os dados atuais):

\begin{Shaded}
\begin{Highlighting}[]
\FunctionTok{sudo}\NormalTok{ docker stop gestoreng{-}web gestoreng{-}worker}
\FunctionTok{sudo}\NormalTok{ docker exec }\AttributeTok{{-}it}\NormalTok{ gestoreng{-}api ./docker{-}entrypoint.sh backup listar}
\FunctionTok{sudo}\NormalTok{ docker exec }\AttributeTok{{-}it}\NormalTok{ gestoreng{-}api ./docker{-}entrypoint.sh backup restaurar gestoreng/gestoreng{-}AAAAMMDD{-}HHMMSS.tar.gz}
\FunctionTok{sudo}\NormalTok{ docker restart gestoreng{-}api }\KeywordTok{\&\&} \FunctionTok{sudo}\NormalTok{ docker start gestoreng{-}worker gestoreng{-}web}
\end{Highlighting}
\end{Shaded}

\hypertarget{problemas-comuns}{%
\section{Problemas comuns}\label{problemas-comuns}}

\textbf{\texttt{invalid\ user\ override:\ generated\ override\ does\ not\ reproduce\ submitted\ Compose}}
--- o ZimaOS não aceita compose com atalhos YAML (\texttt{\&},
\texttt{*}, \texttt{\textless{}\textless{}:}) nem blocos \texttt{x-}
além do \texttt{x-casaos}. Use os arquivos atuais de
\texttt{deploy/zimaos/} (já sem atalhos) e cole o conteúdo inteiro, sem
reformatar.

\textbf{\texttt{lookup\ ghcr.1panel.live:\ no\ such\ host}} --- espelho
de registro inválido no Docker. No 1Panel: \emph{Contêineres
\textgreater{} Configuração \textgreater{} Espelhos de registro}, apague
o espelho e salve. Manualmente:

\begin{Shaded}
\begin{Highlighting}[]
\FunctionTok{sudo}\NormalTok{ cat /etc/docker/daemon.json          }\CommentTok{\# localize "registry{-}mirrors"}
\FunctionTok{sudo}\NormalTok{ nano /etc/docker/daemon.json         }\CommentTok{\# remova o endereço problemático}
\FunctionTok{sudo}\NormalTok{ systemctl restart docker}
\end{Highlighting}
\end{Shaded}

\textbf{\texttt{unauthorized} / \texttt{denied} ao baixar a imagem} ---
faça (ou renove) o login do passo 1.2. Confira se o token tem
\texttt{read:packages} e se sua conta tem acesso ao repositório.

\textbf{\texttt{password\ authentication\ failed\ for\ user}} --- senha
ou usuário diferentes do banco. Confira maiúsculas/minúsculas (passo
4.1) e, no banco interno, lembre que \texttt{POSTGRES\_PASSWORD} só vale
na primeira criação.

\textbf{\texttt{connection\ refused} / \texttt{could\ not\ connect}}
(banco externo) --- refaça os testes do passo 4.2.

\textbf{Pasta de arquivos sem permissão} (aparece no \texttt{verificar}
ou ao anexar arquivos):
\texttt{sudo\ chown\ -R\ 10001:10001\ /media/URSUS/AppData/GestorEng/storage}

\textbf{Sou desconectado ao recarregar a página} ---
\texttt{COOKIE\_SECURE} está \texttt{"true"} mas o acesso é por
\texttt{http://}. Use \texttt{"false"} ou configure HTTPS.

\textbf{Upload de arquivo falha} --- aumente \texttt{MAX\_UPLOAD\_MB} e
\texttt{MAX\_BODY} (iguais).

\textbf{Esqueci a senha do administrador} ---
\texttt{sudo\ docker\ exec\ -it\ gestoreng-api\ python\ -m\ app.manage\ senha\ SEU\_EMAIL}.

\textbf{Worker mostrando ``Aguardando migrações\ldots{}''} --- normal
nos primeiros segundos; se persistir, veja
\texttt{sudo\ docker\ logs\ gestoreng-api}.

\textbf{Porta 8917 já em uso} --- troque por outra livre (veja a tabela
de variáveis).

\textbf{Ver logs de qualquer parte:}

\begin{Shaded}
\begin{Highlighting}[]
\FunctionTok{sudo}\NormalTok{ docker logs }\AttributeTok{{-}{-}tail}\NormalTok{ 100 gestoreng{-}api}
\FunctionTok{sudo}\NormalTok{ docker logs }\AttributeTok{{-}{-}tail}\NormalTok{ 100 gestoreng{-}worker}
\FunctionTok{sudo}\NormalTok{ docker logs }\AttributeTok{{-}{-}tail}\NormalTok{ 100 gestoreng{-}web}
\FunctionTok{sudo}\NormalTok{ docker logs }\AttributeTok{{-}{-}tail}\NormalTok{ 100 gestoreng{-}db      }\CommentTok{\# banco interno}
\end{Highlighting}
\end{Shaded}

\hypertarget{acesso-pelo-celular-e-https-opcional}{%
\section{Acesso pelo celular e HTTPS
(opcional)}\label{acesso-pelo-celular-e-https-opcional}}

\begin{itemize}
\tightlist
\item
  \textbf{App Android:} no primeiro uso informe
  \texttt{http://IP-DO-ZIMAOS:8917} (na rede da empresa) ou o seu
  domínio HTTPS.
\item
  \textbf{HTTPS / acesso externo:} coloque o GestorEng atrás do proxy
  que você já usa (Nginx Proxy Manager, Cloudflare Tunnel, Traefik)
  apontando para \texttt{IP-DO-ZIMAOS:8917}. Depois ajuste
  \texttt{COOKIE\_SECURE:\ "true"} e inclua o domínio em
  \texttt{CORS\_ORIGINS} (ex.: \texttt{https://gestor.velmont.com.br}).
\end{itemize}

\hypertarget{anexo-a-banco-interno.yml}{%
\section{Anexo A ---
banco-interno.yml}\label{anexo-a-banco-interno.yml}}

\begin{quote}
Para instalar, copie do arquivo \texttt{deploy/zimaos/banco-interno.yml}
do repositório: o PDF quebra linhas longas.
\end{quote}

\begin{Shaded}
\begin{Highlighting}[]
\CommentTok{\# GestorEng — ZimaOS · BANCO INTERNO (PostgreSQL + Redis instalados junto com o app)}
\CommentTok{\# Troque todos os valores TROQUE\_*. As variáveis de "api" e "worker" devem ser IGUAIS.}
\CommentTok{\# POSTGRES\_PASSWORD deve ser igual em "db", "api" e "worker"; a senha do Redis igual no "command" do redis e em REDIS\_PASSWORD.}
\CommentTok{\# Senhas podem ter qualquer caractere — o sistema codifica sozinho. Chaves: openssl rand {-}hex 32}
\FunctionTok{name}\KeywordTok{:}\AttributeTok{ gestoreng}
\FunctionTok{services}\KeywordTok{:}
\AttributeTok{  }\FunctionTok{db}\KeywordTok{:}
\AttributeTok{    }\FunctionTok{container\_name}\KeywordTok{:}\AttributeTok{ gestoreng{-}db}
\AttributeTok{    }\FunctionTok{image}\KeywordTok{:}\AttributeTok{ postgres:16{-}alpine}
\AttributeTok{    }\FunctionTok{restart}\KeywordTok{:}\AttributeTok{ unless{-}stopped}
\AttributeTok{    }\FunctionTok{environment}\KeywordTok{:}
\AttributeTok{      }\FunctionTok{POSTGRES\_DB}\KeywordTok{:}\AttributeTok{ gestoreng}
\AttributeTok{      }\FunctionTok{POSTGRES\_USER}\KeywordTok{:}\AttributeTok{ gestoreng}
\AttributeTok{      }\FunctionTok{POSTGRES\_PASSWORD}\KeywordTok{:}\AttributeTok{ TROQUE\_SENHA\_DO\_BANCO}
\AttributeTok{      }\FunctionTok{TZ}\KeywordTok{:}\AttributeTok{ America/Sao\_Paulo}
\AttributeTok{    }\FunctionTok{volumes}\KeywordTok{:}
\AttributeTok{      }\KeywordTok{{-}}\AttributeTok{ }\FunctionTok{type}\KeywordTok{:}\AttributeTok{ bind}
\AttributeTok{        }\FunctionTok{source}\KeywordTok{:}\AttributeTok{ /media/URSUS/AppData/GestorEng/postgres}
\AttributeTok{        }\FunctionTok{target}\KeywordTok{:}\AttributeTok{ /var/lib/postgresql/data}
\AttributeTok{    }\FunctionTok{networks}\KeywordTok{:}
\AttributeTok{      }\KeywordTok{{-}}\AttributeTok{ gestoreng}
\AttributeTok{    }\FunctionTok{security\_opt}\KeywordTok{:}
\AttributeTok{      }\KeywordTok{{-}}\AttributeTok{ no{-}new{-}privileges:true}

\AttributeTok{  }\FunctionTok{redis}\KeywordTok{:}
\AttributeTok{    }\FunctionTok{container\_name}\KeywordTok{:}\AttributeTok{ gestoreng{-}redis}
\AttributeTok{    }\FunctionTok{image}\KeywordTok{:}\AttributeTok{ redis:7{-}alpine}
\AttributeTok{    }\FunctionTok{restart}\KeywordTok{:}\AttributeTok{ unless{-}stopped}
\AttributeTok{    }\FunctionTok{command}\KeywordTok{:}
\AttributeTok{      }\KeywordTok{{-}}\AttributeTok{ redis{-}server}
\AttributeTok{      }\KeywordTok{{-}}\AttributeTok{ {-}{-}requirepass}
\AttributeTok{      }\KeywordTok{{-}}\AttributeTok{ TROQUE\_SENHA\_DO\_REDIS}
\AttributeTok{      }\KeywordTok{{-}}\AttributeTok{ {-}{-}save}
\AttributeTok{      }\KeywordTok{{-}}\AttributeTok{ }\StringTok{""}
\AttributeTok{      }\KeywordTok{{-}}\AttributeTok{ {-}{-}appendonly}
\AttributeTok{      }\KeywordTok{{-}}\AttributeTok{ }\StringTok{"no"}
\AttributeTok{    }\FunctionTok{networks}\KeywordTok{:}
\AttributeTok{      }\KeywordTok{{-}}\AttributeTok{ gestoreng}
\AttributeTok{    }\FunctionTok{read\_only}\KeywordTok{:}\AttributeTok{ }\CharTok{true}
\AttributeTok{    }\FunctionTok{security\_opt}\KeywordTok{:}
\AttributeTok{      }\KeywordTok{{-}}\AttributeTok{ no{-}new{-}privileges:true}

\AttributeTok{  }\FunctionTok{api}\KeywordTok{:}
\AttributeTok{    }\FunctionTok{container\_name}\KeywordTok{:}\AttributeTok{ gestoreng{-}api}
\AttributeTok{    }\FunctionTok{image}\KeywordTok{:}\AttributeTok{ ghcr.io/incorpore/gestoreng{-}api:latest}
\AttributeTok{    }\FunctionTok{restart}\KeywordTok{:}\AttributeTok{ unless{-}stopped}
\AttributeTok{    }\FunctionTok{command}\KeywordTok{:}
\AttributeTok{      }\KeywordTok{{-}}\AttributeTok{ api}
\AttributeTok{    }\FunctionTok{environment}\KeywordTok{:}
\AttributeTok{      }\FunctionTok{ENVIRONMENT}\KeywordTok{:}\AttributeTok{ production}
\AttributeTok{      }\FunctionTok{POSTGRES\_HOST}\KeywordTok{:}\AttributeTok{ db}
\AttributeTok{      }\FunctionTok{POSTGRES\_PORT}\KeywordTok{:}\AttributeTok{ }\StringTok{"5432"}
\AttributeTok{      }\FunctionTok{POSTGRES\_DB}\KeywordTok{:}\AttributeTok{ gestoreng}
\AttributeTok{      }\FunctionTok{POSTGRES\_USER}\KeywordTok{:}\AttributeTok{ gestoreng}
\AttributeTok{      }\FunctionTok{POSTGRES\_PASSWORD}\KeywordTok{:}\AttributeTok{ TROQUE\_SENHA\_DO\_BANCO}
\AttributeTok{      }\FunctionTok{REDIS\_HOST}\KeywordTok{:}\AttributeTok{ redis}
\AttributeTok{      }\FunctionTok{REDIS\_PORT}\KeywordTok{:}\AttributeTok{ }\StringTok{"6379"}
\AttributeTok{      }\FunctionTok{REDIS\_PASSWORD}\KeywordTok{:}\AttributeTok{ TROQUE\_SENHA\_DO\_REDIS}
\AttributeTok{      }\FunctionTok{REDIS\_DB}\KeywordTok{:}\AttributeTok{ }\StringTok{"0"}
\AttributeTok{      }\FunctionTok{REDIS\_PREFIX}\KeywordTok{:}\AttributeTok{ GestorEng\_}
\AttributeTok{      }\FunctionTok{SECRET\_KEY}\KeywordTok{:}\AttributeTok{ TROQUE\_CHAVE\_1}
\AttributeTok{      }\FunctionTok{ENCRYPTION\_KEY}\KeywordTok{:}\AttributeTok{ TROQUE\_CHAVE\_2}
\AttributeTok{      }\FunctionTok{ADMIN\_EMAIL}\KeywordTok{:}\AttributeTok{ admin@velmont.com.br}
\AttributeTok{      }\FunctionTok{ADMIN\_PASSWORD}\KeywordTok{:}\AttributeTok{ TROQUE\_Senha@Forte2026}
\AttributeTok{      }\FunctionTok{ADMIN\_NAME}\KeywordTok{:}\AttributeTok{ Administrador}
\AttributeTok{      }\FunctionTok{COOKIE\_SECURE}\KeywordTok{:}\AttributeTok{ }\StringTok{"false"}
\AttributeTok{      }\FunctionTok{CORS\_ORIGINS}\KeywordTok{:}\AttributeTok{ http://192.168.3.100:8917,https://localhost,capacitor://localhost}
\AttributeTok{      }\FunctionTok{MAX\_UPLOAD\_MB}\KeywordTok{:}\AttributeTok{ }\StringTok{"50"}
\AttributeTok{      }\FunctionTok{TZ}\KeywordTok{:}\AttributeTok{ America/Sao\_Paulo}
\AttributeTok{      }\FunctionTok{WEB\_CONCURRENCY}\KeywordTok{:}\AttributeTok{ }\StringTok{"3"}
\AttributeTok{      }\FunctionTok{STORAGE\_DIR}\KeywordTok{:}\AttributeTok{ /data/storage}
\AttributeTok{    }\FunctionTok{volumes}\KeywordTok{:}
\AttributeTok{      }\KeywordTok{{-}}\AttributeTok{ }\FunctionTok{type}\KeywordTok{:}\AttributeTok{ bind}
\AttributeTok{        }\FunctionTok{source}\KeywordTok{:}\AttributeTok{ /media/URSUS/AppData/GestorEng/storage}
\AttributeTok{        }\FunctionTok{target}\KeywordTok{:}\AttributeTok{ /data/storage}
\AttributeTok{    }\FunctionTok{depends\_on}\KeywordTok{:}
\AttributeTok{      }\KeywordTok{{-}}\AttributeTok{ db}
\AttributeTok{      }\KeywordTok{{-}}\AttributeTok{ redis}
\AttributeTok{    }\FunctionTok{networks}\KeywordTok{:}
\AttributeTok{      }\KeywordTok{{-}}\AttributeTok{ gestoreng}
\AttributeTok{    }\FunctionTok{read\_only}\KeywordTok{:}\AttributeTok{ }\CharTok{true}
\AttributeTok{    }\FunctionTok{tmpfs}\KeywordTok{:}
\AttributeTok{      }\KeywordTok{{-}}\AttributeTok{ /tmp}
\AttributeTok{    }\FunctionTok{security\_opt}\KeywordTok{:}
\AttributeTok{      }\KeywordTok{{-}}\AttributeTok{ no{-}new{-}privileges:true}
\AttributeTok{    }\FunctionTok{cap\_drop}\KeywordTok{:}
\AttributeTok{      }\KeywordTok{{-}}\AttributeTok{ ALL}

\AttributeTok{  }\FunctionTok{worker}\KeywordTok{:}
\AttributeTok{    }\FunctionTok{container\_name}\KeywordTok{:}\AttributeTok{ gestoreng{-}worker}
\AttributeTok{    }\FunctionTok{image}\KeywordTok{:}\AttributeTok{ ghcr.io/incorpore/gestoreng{-}api:latest}
\AttributeTok{    }\FunctionTok{restart}\KeywordTok{:}\AttributeTok{ unless{-}stopped}
\AttributeTok{    }\FunctionTok{command}\KeywordTok{:}
\AttributeTok{      }\KeywordTok{{-}}\AttributeTok{ worker}
\AttributeTok{    }\FunctionTok{environment}\KeywordTok{:}
\AttributeTok{      }\FunctionTok{ENVIRONMENT}\KeywordTok{:}\AttributeTok{ production}
\AttributeTok{      }\FunctionTok{POSTGRES\_HOST}\KeywordTok{:}\AttributeTok{ db}
\AttributeTok{      }\FunctionTok{POSTGRES\_PORT}\KeywordTok{:}\AttributeTok{ }\StringTok{"5432"}
\AttributeTok{      }\FunctionTok{POSTGRES\_DB}\KeywordTok{:}\AttributeTok{ gestoreng}
\AttributeTok{      }\FunctionTok{POSTGRES\_USER}\KeywordTok{:}\AttributeTok{ gestoreng}
\AttributeTok{      }\FunctionTok{POSTGRES\_PASSWORD}\KeywordTok{:}\AttributeTok{ TROQUE\_SENHA\_DO\_BANCO}
\AttributeTok{      }\FunctionTok{REDIS\_HOST}\KeywordTok{:}\AttributeTok{ redis}
\AttributeTok{      }\FunctionTok{REDIS\_PORT}\KeywordTok{:}\AttributeTok{ }\StringTok{"6379"}
\AttributeTok{      }\FunctionTok{REDIS\_PASSWORD}\KeywordTok{:}\AttributeTok{ TROQUE\_SENHA\_DO\_REDIS}
\AttributeTok{      }\FunctionTok{REDIS\_DB}\KeywordTok{:}\AttributeTok{ }\StringTok{"0"}
\AttributeTok{      }\FunctionTok{REDIS\_PREFIX}\KeywordTok{:}\AttributeTok{ GestorEng\_}
\AttributeTok{      }\FunctionTok{SECRET\_KEY}\KeywordTok{:}\AttributeTok{ TROQUE\_CHAVE\_1}
\AttributeTok{      }\FunctionTok{ENCRYPTION\_KEY}\KeywordTok{:}\AttributeTok{ TROQUE\_CHAVE\_2}
\AttributeTok{      }\FunctionTok{COOKIE\_SECURE}\KeywordTok{:}\AttributeTok{ }\StringTok{"false"}
\AttributeTok{      }\FunctionTok{CORS\_ORIGINS}\KeywordTok{:}\AttributeTok{ http://192.168.3.100:8917,https://localhost,capacitor://localhost}
\AttributeTok{      }\FunctionTok{MAX\_UPLOAD\_MB}\KeywordTok{:}\AttributeTok{ }\StringTok{"50"}
\AttributeTok{      }\FunctionTok{TZ}\KeywordTok{:}\AttributeTok{ America/Sao\_Paulo}
\AttributeTok{      }\FunctionTok{STORAGE\_DIR}\KeywordTok{:}\AttributeTok{ /data/storage}
\AttributeTok{    }\FunctionTok{volumes}\KeywordTok{:}
\AttributeTok{      }\KeywordTok{{-}}\AttributeTok{ }\FunctionTok{type}\KeywordTok{:}\AttributeTok{ bind}
\AttributeTok{        }\FunctionTok{source}\KeywordTok{:}\AttributeTok{ /media/URSUS/AppData/GestorEng/storage}
\AttributeTok{        }\FunctionTok{target}\KeywordTok{:}\AttributeTok{ /data/storage}
\AttributeTok{    }\FunctionTok{depends\_on}\KeywordTok{:}
\AttributeTok{      }\KeywordTok{{-}}\AttributeTok{ api}
\AttributeTok{    }\FunctionTok{healthcheck}\KeywordTok{:}
\AttributeTok{      }\FunctionTok{disable}\KeywordTok{:}\AttributeTok{ }\CharTok{true}
\AttributeTok{    }\FunctionTok{networks}\KeywordTok{:}
\AttributeTok{      }\KeywordTok{{-}}\AttributeTok{ gestoreng}
\AttributeTok{    }\FunctionTok{read\_only}\KeywordTok{:}\AttributeTok{ }\CharTok{true}
\AttributeTok{    }\FunctionTok{tmpfs}\KeywordTok{:}
\AttributeTok{      }\KeywordTok{{-}}\AttributeTok{ /tmp}
\AttributeTok{    }\FunctionTok{security\_opt}\KeywordTok{:}
\AttributeTok{      }\KeywordTok{{-}}\AttributeTok{ no{-}new{-}privileges:true}
\AttributeTok{    }\FunctionTok{cap\_drop}\KeywordTok{:}
\AttributeTok{      }\KeywordTok{{-}}\AttributeTok{ ALL}

\AttributeTok{  }\FunctionTok{web}\KeywordTok{:}
\AttributeTok{    }\FunctionTok{container\_name}\KeywordTok{:}\AttributeTok{ gestoreng{-}web}
\AttributeTok{    }\FunctionTok{image}\KeywordTok{:}\AttributeTok{ ghcr.io/incorpore/gestoreng{-}web:latest}
\AttributeTok{    }\FunctionTok{restart}\KeywordTok{:}\AttributeTok{ unless{-}stopped}
\AttributeTok{    }\FunctionTok{environment}\KeywordTok{:}
\AttributeTok{      }\FunctionTok{API\_UPSTREAM}\KeywordTok{:}\AttributeTok{ api:8000}
\AttributeTok{      }\FunctionTok{MAX\_BODY}\KeywordTok{:}\AttributeTok{ 50m}
\AttributeTok{    }\FunctionTok{ports}\KeywordTok{:}
\AttributeTok{      }\KeywordTok{{-}}\AttributeTok{ }\FunctionTok{target}\KeywordTok{:}\AttributeTok{ }\DecValTok{8080}
\AttributeTok{        }\FunctionTok{published}\KeywordTok{:}\AttributeTok{ }\StringTok{"8917"}
\AttributeTok{        }\FunctionTok{protocol}\KeywordTok{:}\AttributeTok{ tcp}
\AttributeTok{    }\FunctionTok{depends\_on}\KeywordTok{:}
\AttributeTok{      }\KeywordTok{{-}}\AttributeTok{ api}
\AttributeTok{    }\FunctionTok{networks}\KeywordTok{:}
\AttributeTok{      }\KeywordTok{{-}}\AttributeTok{ gestoreng}
\AttributeTok{    }\FunctionTok{read\_only}\KeywordTok{:}\AttributeTok{ }\CharTok{true}
\AttributeTok{    }\FunctionTok{tmpfs}\KeywordTok{:}
\AttributeTok{      }\KeywordTok{{-}}\AttributeTok{ /tmp}
\AttributeTok{      }\KeywordTok{{-}}\AttributeTok{ /etc/nginx/conf.d}
\AttributeTok{      }\KeywordTok{{-}}\AttributeTok{ /var/cache/nginx}
\AttributeTok{    }\FunctionTok{security\_opt}\KeywordTok{:}
\AttributeTok{      }\KeywordTok{{-}}\AttributeTok{ no{-}new{-}privileges:true}
\AttributeTok{    }\FunctionTok{cap\_drop}\KeywordTok{:}
\AttributeTok{      }\KeywordTok{{-}}\AttributeTok{ ALL}

\FunctionTok{networks}\KeywordTok{:}
\AttributeTok{  }\FunctionTok{gestoreng}\KeywordTok{:}
\AttributeTok{    }\FunctionTok{name}\KeywordTok{:}\AttributeTok{ gestoreng}
\AttributeTok{    }\FunctionTok{driver}\KeywordTok{:}\AttributeTok{ bridge}

\FunctionTok{x{-}casaos}\KeywordTok{:}
\AttributeTok{  }\FunctionTok{author}\KeywordTok{:}\AttributeTok{ Velmont}
\AttributeTok{  }\FunctionTok{category}\KeywordTok{:}\AttributeTok{ Productivity}
\AttributeTok{  }\FunctionTok{description}\KeywordTok{:}
\AttributeTok{    }\FunctionTok{custom}\KeywordTok{:}\AttributeTok{ Sistema de Gestão de Engenharia — obras, licitações, financeiro e documentos}
\AttributeTok{  }\FunctionTok{index}\KeywordTok{:}\AttributeTok{ /}
\AttributeTok{  }\FunctionTok{main}\KeywordTok{:}\AttributeTok{ web}
\AttributeTok{  }\FunctionTok{port\_map}\KeywordTok{:}\AttributeTok{ }\StringTok{"8917"}
\AttributeTok{  }\FunctionTok{scheme}\KeywordTok{:}\AttributeTok{ http}
\AttributeTok{  }\FunctionTok{title}\KeywordTok{:}
\AttributeTok{    }\FunctionTok{custom}\KeywordTok{:}\AttributeTok{ GestorEng}
\end{Highlighting}
\end{Shaded}

\hypertarget{anexo-b-banco-externo.yml}{%
\section{Anexo B ---
banco-externo.yml}\label{anexo-b-banco-externo.yml}}

\begin{quote}
Para instalar, copie do arquivo \texttt{deploy/zimaos/banco-externo.yml}
do repositório: o PDF quebra linhas longas.
\end{quote}

\begin{Shaded}
\begin{Highlighting}[]
\CommentTok{\# GestorEng — ZimaOS · BANCO EXTERNO (PostgreSQL e Redis já existentes na rede)}
\CommentTok{\# Troque todos os valores TROQUE\_*. As variáveis de "api" e "worker" devem ser IGUAIS.}
\CommentTok{\# Redis é opcional: para não usar, apague as linhas REDIS\_* de "api" e "worker".}
\CommentTok{\# Senhas podem ter qualquer caractere — o sistema codifica sozinho. Chaves: openssl rand {-}hex 32}
\FunctionTok{name}\KeywordTok{:}\AttributeTok{ gestoreng}
\FunctionTok{services}\KeywordTok{:}
\AttributeTok{  }\FunctionTok{api}\KeywordTok{:}
\AttributeTok{    }\FunctionTok{container\_name}\KeywordTok{:}\AttributeTok{ gestoreng{-}api}
\AttributeTok{    }\FunctionTok{image}\KeywordTok{:}\AttributeTok{ ghcr.io/incorpore/gestoreng{-}api:latest}
\AttributeTok{    }\FunctionTok{restart}\KeywordTok{:}\AttributeTok{ unless{-}stopped}
\AttributeTok{    }\FunctionTok{command}\KeywordTok{:}
\AttributeTok{      }\KeywordTok{{-}}\AttributeTok{ api}
\AttributeTok{    }\FunctionTok{environment}\KeywordTok{:}
\AttributeTok{      }\FunctionTok{ENVIRONMENT}\KeywordTok{:}\AttributeTok{ production}
\AttributeTok{      }\FunctionTok{POSTGRES\_HOST}\KeywordTok{:}\AttributeTok{ }\FloatTok{192.168.3.100}
\AttributeTok{      }\FunctionTok{POSTGRES\_PORT}\KeywordTok{:}\AttributeTok{ }\StringTok{"5432"}
\AttributeTok{      }\FunctionTok{POSTGRES\_DB}\KeywordTok{:}\AttributeTok{ gestoreng}
\AttributeTok{      }\FunctionTok{POSTGRES\_USER}\KeywordTok{:}\AttributeTok{ gestoreng}
\AttributeTok{      }\FunctionTok{POSTGRES\_PASSWORD}\KeywordTok{:}\AttributeTok{ TROQUE\_SENHA\_DO\_BANCO}
\AttributeTok{      }\FunctionTok{REDIS\_HOST}\KeywordTok{:}\AttributeTok{ }\FloatTok{192.168.3.100}
\AttributeTok{      }\FunctionTok{REDIS\_PORT}\KeywordTok{:}\AttributeTok{ }\StringTok{"6379"}
\AttributeTok{      }\FunctionTok{REDIS\_PASSWORD}\KeywordTok{:}\AttributeTok{ TROQUE\_SENHA\_DO\_REDIS}
\AttributeTok{      }\FunctionTok{REDIS\_DB}\KeywordTok{:}\AttributeTok{ }\StringTok{"0"}
\AttributeTok{      }\FunctionTok{REDIS\_PREFIX}\KeywordTok{:}\AttributeTok{ GestorEng\_}
\AttributeTok{      }\FunctionTok{SECRET\_KEY}\KeywordTok{:}\AttributeTok{ TROQUE\_CHAVE\_1}
\AttributeTok{      }\FunctionTok{ENCRYPTION\_KEY}\KeywordTok{:}\AttributeTok{ TROQUE\_CHAVE\_2}
\AttributeTok{      }\FunctionTok{ADMIN\_EMAIL}\KeywordTok{:}\AttributeTok{ admin@velmont.com.br}
\AttributeTok{      }\FunctionTok{ADMIN\_PASSWORD}\KeywordTok{:}\AttributeTok{ TROQUE\_Senha@Forte2026}
\AttributeTok{      }\FunctionTok{ADMIN\_NAME}\KeywordTok{:}\AttributeTok{ Administrador}
\AttributeTok{      }\FunctionTok{COOKIE\_SECURE}\KeywordTok{:}\AttributeTok{ }\StringTok{"false"}
\AttributeTok{      }\FunctionTok{CORS\_ORIGINS}\KeywordTok{:}\AttributeTok{ http://192.168.3.100:8917,https://localhost,capacitor://localhost}
\AttributeTok{      }\FunctionTok{MAX\_UPLOAD\_MB}\KeywordTok{:}\AttributeTok{ }\StringTok{"50"}
\AttributeTok{      }\FunctionTok{TZ}\KeywordTok{:}\AttributeTok{ America/Sao\_Paulo}
\AttributeTok{      }\FunctionTok{WEB\_CONCURRENCY}\KeywordTok{:}\AttributeTok{ }\StringTok{"3"}
\AttributeTok{      }\FunctionTok{STORAGE\_DIR}\KeywordTok{:}\AttributeTok{ /data/storage}
\AttributeTok{    }\FunctionTok{volumes}\KeywordTok{:}
\AttributeTok{      }\KeywordTok{{-}}\AttributeTok{ }\FunctionTok{type}\KeywordTok{:}\AttributeTok{ bind}
\AttributeTok{        }\FunctionTok{source}\KeywordTok{:}\AttributeTok{ /media/URSUS/AppData/GestorEng/storage}
\AttributeTok{        }\FunctionTok{target}\KeywordTok{:}\AttributeTok{ /data/storage}
\AttributeTok{    }\FunctionTok{networks}\KeywordTok{:}
\AttributeTok{      }\KeywordTok{{-}}\AttributeTok{ gestoreng}
\AttributeTok{    }\FunctionTok{read\_only}\KeywordTok{:}\AttributeTok{ }\CharTok{true}
\AttributeTok{    }\FunctionTok{tmpfs}\KeywordTok{:}
\AttributeTok{      }\KeywordTok{{-}}\AttributeTok{ /tmp}
\AttributeTok{    }\FunctionTok{security\_opt}\KeywordTok{:}
\AttributeTok{      }\KeywordTok{{-}}\AttributeTok{ no{-}new{-}privileges:true}
\AttributeTok{    }\FunctionTok{cap\_drop}\KeywordTok{:}
\AttributeTok{      }\KeywordTok{{-}}\AttributeTok{ ALL}

\AttributeTok{  }\FunctionTok{worker}\KeywordTok{:}
\AttributeTok{    }\FunctionTok{container\_name}\KeywordTok{:}\AttributeTok{ gestoreng{-}worker}
\AttributeTok{    }\FunctionTok{image}\KeywordTok{:}\AttributeTok{ ghcr.io/incorpore/gestoreng{-}api:latest}
\AttributeTok{    }\FunctionTok{restart}\KeywordTok{:}\AttributeTok{ unless{-}stopped}
\AttributeTok{    }\FunctionTok{command}\KeywordTok{:}
\AttributeTok{      }\KeywordTok{{-}}\AttributeTok{ worker}
\AttributeTok{    }\FunctionTok{environment}\KeywordTok{:}
\AttributeTok{      }\FunctionTok{ENVIRONMENT}\KeywordTok{:}\AttributeTok{ production}
\AttributeTok{      }\FunctionTok{POSTGRES\_HOST}\KeywordTok{:}\AttributeTok{ }\FloatTok{192.168.3.100}
\AttributeTok{      }\FunctionTok{POSTGRES\_PORT}\KeywordTok{:}\AttributeTok{ }\StringTok{"5432"}
\AttributeTok{      }\FunctionTok{POSTGRES\_DB}\KeywordTok{:}\AttributeTok{ gestoreng}
\AttributeTok{      }\FunctionTok{POSTGRES\_USER}\KeywordTok{:}\AttributeTok{ gestoreng}
\AttributeTok{      }\FunctionTok{POSTGRES\_PASSWORD}\KeywordTok{:}\AttributeTok{ TROQUE\_SENHA\_DO\_BANCO}
\AttributeTok{      }\FunctionTok{REDIS\_HOST}\KeywordTok{:}\AttributeTok{ }\FloatTok{192.168.3.100}
\AttributeTok{      }\FunctionTok{REDIS\_PORT}\KeywordTok{:}\AttributeTok{ }\StringTok{"6379"}
\AttributeTok{      }\FunctionTok{REDIS\_PASSWORD}\KeywordTok{:}\AttributeTok{ TROQUE\_SENHA\_DO\_REDIS}
\AttributeTok{      }\FunctionTok{REDIS\_DB}\KeywordTok{:}\AttributeTok{ }\StringTok{"0"}
\AttributeTok{      }\FunctionTok{REDIS\_PREFIX}\KeywordTok{:}\AttributeTok{ GestorEng\_}
\AttributeTok{      }\FunctionTok{SECRET\_KEY}\KeywordTok{:}\AttributeTok{ TROQUE\_CHAVE\_1}
\AttributeTok{      }\FunctionTok{ENCRYPTION\_KEY}\KeywordTok{:}\AttributeTok{ TROQUE\_CHAVE\_2}
\AttributeTok{      }\FunctionTok{COOKIE\_SECURE}\KeywordTok{:}\AttributeTok{ }\StringTok{"false"}
\AttributeTok{      }\FunctionTok{CORS\_ORIGINS}\KeywordTok{:}\AttributeTok{ http://192.168.3.100:8917,https://localhost,capacitor://localhost}
\AttributeTok{      }\FunctionTok{MAX\_UPLOAD\_MB}\KeywordTok{:}\AttributeTok{ }\StringTok{"50"}
\AttributeTok{      }\FunctionTok{TZ}\KeywordTok{:}\AttributeTok{ America/Sao\_Paulo}
\AttributeTok{      }\FunctionTok{STORAGE\_DIR}\KeywordTok{:}\AttributeTok{ /data/storage}
\AttributeTok{    }\FunctionTok{volumes}\KeywordTok{:}
\AttributeTok{      }\KeywordTok{{-}}\AttributeTok{ }\FunctionTok{type}\KeywordTok{:}\AttributeTok{ bind}
\AttributeTok{        }\FunctionTok{source}\KeywordTok{:}\AttributeTok{ /media/URSUS/AppData/GestorEng/storage}
\AttributeTok{        }\FunctionTok{target}\KeywordTok{:}\AttributeTok{ /data/storage}
\AttributeTok{    }\FunctionTok{depends\_on}\KeywordTok{:}
\AttributeTok{      }\KeywordTok{{-}}\AttributeTok{ api}
\AttributeTok{    }\FunctionTok{healthcheck}\KeywordTok{:}
\AttributeTok{      }\FunctionTok{disable}\KeywordTok{:}\AttributeTok{ }\CharTok{true}
\AttributeTok{    }\FunctionTok{networks}\KeywordTok{:}
\AttributeTok{      }\KeywordTok{{-}}\AttributeTok{ gestoreng}
\AttributeTok{    }\FunctionTok{read\_only}\KeywordTok{:}\AttributeTok{ }\CharTok{true}
\AttributeTok{    }\FunctionTok{tmpfs}\KeywordTok{:}
\AttributeTok{      }\KeywordTok{{-}}\AttributeTok{ /tmp}
\AttributeTok{    }\FunctionTok{security\_opt}\KeywordTok{:}
\AttributeTok{      }\KeywordTok{{-}}\AttributeTok{ no{-}new{-}privileges:true}
\AttributeTok{    }\FunctionTok{cap\_drop}\KeywordTok{:}
\AttributeTok{      }\KeywordTok{{-}}\AttributeTok{ ALL}

\AttributeTok{  }\FunctionTok{web}\KeywordTok{:}
\AttributeTok{    }\FunctionTok{container\_name}\KeywordTok{:}\AttributeTok{ gestoreng{-}web}
\AttributeTok{    }\FunctionTok{image}\KeywordTok{:}\AttributeTok{ ghcr.io/incorpore/gestoreng{-}web:latest}
\AttributeTok{    }\FunctionTok{restart}\KeywordTok{:}\AttributeTok{ unless{-}stopped}
\AttributeTok{    }\FunctionTok{environment}\KeywordTok{:}
\AttributeTok{      }\FunctionTok{API\_UPSTREAM}\KeywordTok{:}\AttributeTok{ api:8000}
\AttributeTok{      }\FunctionTok{MAX\_BODY}\KeywordTok{:}\AttributeTok{ 50m}
\AttributeTok{    }\FunctionTok{ports}\KeywordTok{:}
\AttributeTok{      }\KeywordTok{{-}}\AttributeTok{ }\FunctionTok{target}\KeywordTok{:}\AttributeTok{ }\DecValTok{8080}
\AttributeTok{        }\FunctionTok{published}\KeywordTok{:}\AttributeTok{ }\StringTok{"8917"}
\AttributeTok{        }\FunctionTok{protocol}\KeywordTok{:}\AttributeTok{ tcp}
\AttributeTok{    }\FunctionTok{depends\_on}\KeywordTok{:}
\AttributeTok{      }\KeywordTok{{-}}\AttributeTok{ api}
\AttributeTok{    }\FunctionTok{networks}\KeywordTok{:}
\AttributeTok{      }\KeywordTok{{-}}\AttributeTok{ gestoreng}
\AttributeTok{    }\FunctionTok{read\_only}\KeywordTok{:}\AttributeTok{ }\CharTok{true}
\AttributeTok{    }\FunctionTok{tmpfs}\KeywordTok{:}
\AttributeTok{      }\KeywordTok{{-}}\AttributeTok{ /tmp}
\AttributeTok{      }\KeywordTok{{-}}\AttributeTok{ /etc/nginx/conf.d}
\AttributeTok{      }\KeywordTok{{-}}\AttributeTok{ /var/cache/nginx}
\AttributeTok{    }\FunctionTok{security\_opt}\KeywordTok{:}
\AttributeTok{      }\KeywordTok{{-}}\AttributeTok{ no{-}new{-}privileges:true}
\AttributeTok{    }\FunctionTok{cap\_drop}\KeywordTok{:}
\AttributeTok{      }\KeywordTok{{-}}\AttributeTok{ ALL}

\FunctionTok{networks}\KeywordTok{:}
\AttributeTok{  }\FunctionTok{gestoreng}\KeywordTok{:}
\AttributeTok{    }\FunctionTok{name}\KeywordTok{:}\AttributeTok{ gestoreng}
\AttributeTok{    }\FunctionTok{driver}\KeywordTok{:}\AttributeTok{ bridge}

\FunctionTok{x{-}casaos}\KeywordTok{:}
\AttributeTok{  }\FunctionTok{author}\KeywordTok{:}\AttributeTok{ Velmont}
\AttributeTok{  }\FunctionTok{category}\KeywordTok{:}\AttributeTok{ Productivity}
\AttributeTok{  }\FunctionTok{description}\KeywordTok{:}
\AttributeTok{    }\FunctionTok{custom}\KeywordTok{:}\AttributeTok{ Sistema de Gestão de Engenharia — obras, licitações, financeiro e documentos}
\AttributeTok{  }\FunctionTok{index}\KeywordTok{:}\AttributeTok{ /}
\AttributeTok{  }\FunctionTok{main}\KeywordTok{:}\AttributeTok{ web}
\AttributeTok{  }\FunctionTok{port\_map}\KeywordTok{:}\AttributeTok{ }\StringTok{"8917"}
\AttributeTok{  }\FunctionTok{scheme}\KeywordTok{:}\AttributeTok{ http}
\AttributeTok{  }\FunctionTok{title}\KeywordTok{:}
\AttributeTok{    }\FunctionTok{custom}\KeywordTok{:}\AttributeTok{ GestorEng}
\end{Highlighting}
\end{Shaded}


\end{document}
